
\def\PUDcodigo{EME104}

\def\PUDcodacad{04506.55}

\def\PUDdisciplina{Gestão de Projetos}

\def\PUDsemestre{Opt}

\def\PUDhoras{80}

%NOVOS ---------------------------------------------------
\def\PUDhorasTeoria{}
\def\PUDhorasPratica{}

\def\PUDinterECA{---}

\def\PUDatuaisECA{---}

\def\PUDrecursos{Projetor multimídia; Quadro branco e pincel.}

%----------------------------------------------------------

\def\PUDcreditos{4}

\def\PUDprerequisito{---}

\def\PUDpreacad{---}

\def\PUDobjetivos{Ser capar de gerenciar projetos de engenharia; Conhecer ferramentas de gestão de projetos.}

\def\PUDementa{Estratégia e Projetos;  Estrutura e Etapas do Projeto;  Seleção de Projetos;  Análise Econômica de Projetos;  Gerenciamento de Projetos Seguindo o PMBOK.}

\def\PUDconteudo{ & Estratégias e projetos. 
\\ 
 & Estratégia da organização e seleção de projetos. 
\\ 
 & Estruturas e etapas de um projeto. 
\\ 
 & Análise econômica de projetos. 
\\ 
 & Impactos sociais e ambientais ocasionados pelos projetos. 
\\ 
 & Gerenciamento de projetos seguindo o PMBOK. 
\\ 
 & Processos de gerenciamento de projetos: iniciação, planejamento, execução, monitoramento, controle e encerramento. 
\\ 
 & Áreas de conhecimento do gerenciamento de projetos: integração do projeto, escopo do projeto, tempo do projeto, custos associados ao projeto, qualidade do projeto, recursos humanos, comunicações, riscos do projeto e aquisições. }

\def\PUDmetodo{Aulas expositórias que podem ser teóricas e/ou práticas, onde as práticas no laboratório serão marcadas no decorrer da disciplina.}

\def\PUDavalia{A avaliação será feita com aplicação de provas teóricas e/ou práticas, além da possibilidade de inclusão de trabalhos e seminários no decorrer da disciplina.}

\def\PUDbibbasica{ &  \href{www.biblioteca.ifce.edu.br/index.asp?codigo_sophia=61750}{Maximiano}, Antonio Cesar Amaru.  Administração para empreendedores.  2º ed.  São Paulo, SP: Pearson Prentice Hall, 2011.  240p.  ISBN: 9788576058762. 
\\ 
 & \href{www.biblioteca.ifce.edu.br/index.asp?codigo_sophia=65472}{Anunciação}, Heverton.  Gestão de projetos nas melhores práticas para satisfazer o consumidor 2.0.  1º ed.  Rio de Janeiro - RJ.  Editora Ciência Moderna.  2009.  354p.  ISBN: 9788573937787.
% &   GIDO, J.;  CLEMENTS, J.  P.;  Gestão de Projetos.  1ª edição.  São Paulo.  Editora CENGAGE.  2007.  452p.  ISBN 9788522105557. 
\\ 
 %&  Oliveira, Guilherme Bueno de.  Microsoft Project 2010 e Gestão de Projetos.  E-book. Editora Pearson.  306p.  ISBN: 9788576059523. 
 
 &  \href{www.biblioteca.ifce.edu.br/index.asp?codigo_sophia=40465}{Pahl}, Gerhard.  Projeto na engenharia: fundamentos do desenvolvimento eficaz de produtos, métodos e aplicaçõe.  São Paulo, SP: Blucher, 2013.  412p.  ISBN: 9788521203636. }

\def\PUDbibcomplementar{ &  \href{www.biblioteca.ifce.edu.br/index.asp?codigo_sophia=61801}{Molinari}, Leonardo.  Gestão de Projetos - Teoria, Técnicas e Práticas.  1º ed.  São Paulo.  Editora Érica.  2010.  240p.  ISBN: 9788536502762.
\\ 
 &  \href{www.biblioteca.ifce.edu.br/index.asp?codigo_sophia=61777}{Newton}, R.  O gestor de Projetos.  2º ed.  São Paulo.  Editora Pearson BRASIL.  2011.  320p.  ISBN: 9788576058113.  
\\ 
 &  \href{www.biblioteca.ifce.edu.br/index.asp?codigo_sophia=65145}{Madureira}, O.  M.  Metodologia do Projeto: planejamento, execução e gerenciamento.  2º ed.  São Paulo.  Editora Blucher.  2015. 361p. ISBN: 9788521209133.
\\
 &  \href{www.biblioteca.ifce.edu.br/index.asp?codigo_sophia=56916}{Marras}, Jean Pierre.  Administração de remuneração. E-book.  Pearson.  244p.  ISBN: 9788581430904.
\\
 &  \href{www.	biblioteca.ifce.edu.br/index.asp?codigo_sophia=55454}{Martelanc}, Roy; Pasin, Rodrigo; Cavalcante, Francisco.  Avaliação de Empresas: um guia para fusões e aquisições e gestão de valor.  E-book.  Pearson.  302p.  ISBN: 9788576050087. }

\def\PUDautor{Venceslau Xavier de Lima Filho}

\def\PUDdata{2013-06-16}

\def\PUDrevisao{1}

\def\PUDrevisor{Samuel Vieira Dias}

\def\PUDdatarevisao{2017-04-30}

\PUDcorpo
